\section{Conclusion}
This thesis set out to investigate how data-layout correctness could be
approached through language design rather than through tooling
or external verification. It proposed extensions to \textit{Cedar},
a declarative language and compiler that allows programmers to specify
physical representation of data structures in C without explicit, byte-level
precision. The project asked whether such a language could
express realistic layouts while remaining compatible with existing
C toolchains and conventional programming practices.

The research demonstrated that it can. Cedar's implementation
shows that layout intent, often buried in compiler pragmas
or implicit alignment rules, can be expressed directly
and compiled into deterministic, verified C code.
Through a compact syntax based on placement directives and compositional
reasoning about representations as a first-class concern. Its compiler
pipeline, structured through intermediate representations, provides
a transparent mapping from high-level specification to concrete bit
ranges. The evaluation confirmed that this approach
faithfully reproduces real-world layouts, including
nested structures and externally defined formats
such as the ELF header, with results consistent across compilers.

Beyond empirical correctness, Cedar contributes
conceptually by reframing layout design as a domain of reasoning.
The language demonstrates that explicitness alone can
eliminate many of the ambiguities that
make low-level code fragile. It treats data representation
not as a by-product of compilation, but as a structured
subject to analysis and automation. This offers
an alternative perspective on systems programming: one in which
correctness emerges from specification rather than from
defensive debugging.

At the same time, the work highlights clear opportunities for a continuation.
Completing semantic checking, extending type coverage, and integrating
with verified backends such as CompCert would advance Cedar as a tool
in an ecosystem of full assurance. These directions promise not only a
more powerful tool but also a research vehicle for studying how
declarative specifications interact with formal proofs and real-world
compilers.

In summary, Cedar represents a step towards more reliable and trustworthy
systems programming. It bridges the clarity of declarative design, with the
pragmatism and industry relevance of C, showing that many of the
benefits can be achieved through language structure alone. The
project stands as both a functioning compiler and a methodological argument:
that improving how we \textit{describe} data is the first step
towards improving how we \textit{trust} it.